\chapter{Metoderørsla og avvisinga}

\begin{motto}
Det fanst eit tiår då faget bygde seg ein metode.\\ Det merkelegaste er ikkje at han ikkje heldt.\\ Det er at dei som bygde han, var dei fyrste til å gå.
\end{motto}

Imperial College, South Kensington, september 1962. Nokre titals menn i jakke og slips, med notatblokker og piper, samla i tre dagar under ein tittel som åleine er verd å stogge ved: \textit{Conference on Systematic and Intuitive Methods in Engineering, Industrial Design, Architecture and Communications}\autocite{jones1963conference}. Dei to orda «systematic» og «intuitive» står ved sida av kvarandre i tittelen fordi forsamlinga enno håpa dei kunne forsonast. Vert J.\,Christopher Jones, ingeniør hjå AEI i Manchester, og Denis Thornley frå arkitekturskulen same stad, hadde kalla saman folk som for det meste hadde lært handverket sitt under krigen, ved å plassere radarstasjonar og optimere konvoiruter. Dei såg ingen grunn til at den same klare tenkinga ikkje skulle kunne plassere eit dørhandtak eller ein bydel. I salen vert det snakka om flytskjema, om morfologiske boksar, om å bryte eit problem ned i parametrar og setje dei saman att. Ingen snakkar om geni, om inspirasjon, om det uutsegnelege.

Det dei var ute etter, var å hente forma ut av det private. Så lenge forma «kjem» til meisteren i ein indre prosess som ingen kan opne, kan ho ikkje lærast ut anna enn ved å stå ved meisterens side og herme, ho kan ikkje grunngjevast anna enn med «slik kjenner eg det rett», og ho kan ikkje granskast utanfrå. Metoden skulle bryte denne lukka: dele prosessen i steg som kunne namngjevast, skrivast ned, lærast bort og prøvast. Det var den same trua som dreiv Ulm, no sett i system. To år før konferansen hadde ein arkitekt på kring tjuefem, Christopher Alexander, levert rørsla si mest ambisiøse bok, \textit{Notes on the Synthesis of Form} (1964), der han med mengdelære og grafteori prøvde å vise korleis eit designproblem kunne brytast ned i delproblem og setjast saman att til ei form som passa krava\autocite{alexander1964notes}. I 1970 kom verket som gav heile rørsla namnet sitt, Jones si \textit{Design Methods: Seeds of Human Futures}\autocite{jones1970methods}. På under eit tiår hadde faget skaffa seg konferansar, bøker, eit ordforråd og ei kjensle av at det endeleg var i ferd med å vekse opp. Aktørane i rommet samanlikna seg med ingeniørane dei kom frå, og trudde dei var på veg til same slag fag. Dei tok feil om kva slag fag det var, men i 1962 høyrde ingen den komande tvilen.

\section{Når byggmeistrane gjekk}

For så skjer det som gjer dette kapitlet ulikt alt anna i designhistoria. I 1971, sju år etter \textit{Notes} og medan rørsla framleis var i framgang, skreiv Alexander i \textit{DMG Newsletter}, organet til rørsla sjølv, ein tekst om tilstanden i designmetodefeltet der han tok offentleg avstand frå si eiga bok\autocite{alexander1971state}. Den boka som meir enn nokon hadde gjeve metoderørsla matematisk autoritet, knapt orka han å opne. Diagramma og grafane han ein gong var stolt av, såg han no på som ei lærd avsporing frå det einaste som tel: å lage ting som er gode å leve med. Til folk som studerte designmetode i staden for å designe, hadde han eitt råd, og det har gått inn i forteljinga om faget: gløym heile greia, og gå i gang med å byggje.

Det er ikkje ein utanforståande kritikar som peikar på veikskapar. Det er hovudarkitekten som går ut av sitt eige byggverk medan det enno står, og ropar til dei som er att at dei kastar bort livet. Og han var ikkje åleine. Jones, som arrangerte den fyrste konferansen og skreiv boka som gav rørsla namnet, dreiv i åra etter bort frå den rasjonalistiske kjernen i si eiga bok, mot ein opnare, meir leikande praksis, og kom til å mistru sjølve draumen om å mekanisere designprosessen. Mannen som døypte rørsla, mistrudde til slutt det namnet stod for. To menn, men ikkje to innfall: eit mønster. Dei som gjekk djupast inn i metoden, og difor såg klårast kva han kunne og ikkje kunne, var dei som forlét honom fyrst og hardast. Dei som heldt fast, arbeidde stort sett lenger ute, der prosedyren enno kunne sjå heil ut fordi ingen hadde ført honom heilt til botnen av eit verkeleg designval.

Tenk om Darwin i 1865 hadde kalla evolusjonstanken ei avsporing, eller grunnleggjarane av termodynamikken teke avstand frå hovudsetningane sine eit tiår etter. Vi ville slutta, med rette, at fundamentet aldri heldt. Alexander og Jones forlét ikkje metoden fordi dei var lunefulle. Dei oppdaga at den eksplisitte prosedyren ikkje fanga det som faktisk skjer når ei god form vert til, og at å dyrke prosedyren gjorde tinga verre. Men dei kunne ikkje seie kvifor på ein måte som gav ein \emph{betre} metode. Dei forkasta honom og gjekk attende til intuisjonen, til atelieret. Rørsla enda såleis ikkje med ein korrigert metode, slik ein vitskap korrigerer ein hypotese. Ho enda med ein retrett frå sjølve tanken om metode.

\section{Det tomme valsteget}

Kvifor greidde faget ikkje halde på ein metode, sjølv med så kloke hovud som Alexander og Jones? Ein metode er ikkje sjølvberande. Han er ein framgangsmåte for å kome fram til noko, og han føreset at ein veit kva det noko er. Den vitskaplege metoden kviler på ein teori om kva kunnskap og kva eit prov er. Eit moderne fag ber metodane sine på eit fundament under seg som fortel kva «rett» tyder før framgangsmåten set i gang. Metoderørsla bygde ein framgangsmåte for å kome fram til den rette forma, men under han fanst det ingen teori om form som kunne seie kva «rett form» var. Dette er ikkje å seie at design er ein ingeniørdisiplin som ikkje fekk det til. Design er noko anna slag, og må grunngjevast på sine eigne premiss. Poenget er smalare og hardare: metoden lova å fjerne den uutsegnelege dommen, og fjerna han ikkje. Stega kunne ein skrive: samle krav, generer alternativ, vurder. Men i det eine steget som avgjer alt, der eitt alternativ skal kårast som betre enn eit anna, hadde metoden ingenting å stø seg til utan nett den smaken og det auget han var meint å avløyse. Han skuva den uutsegnelege dommen eitt steg lenger inn i prosessen. Han hadde ingenting å setje i staden.

\begin{kasus}{Nedbrytinga som stoppa ved terskelen}
Det opphavlege dømet i \textit{Notes} er eit indisk landsbyplan-problem. Alexander braut det ned i over hundre krav, det boka kallar «misfits», og kopla dei i eit nett av diagram. Maskineriet var imponerande: motsetnader mellom krav vart kartlagde, klyngjer av samanhengande krav vart funne, og kvar klynge skulle svare til ein «del» av forma. Men ved det avgjerande punktet, der ein klynge av krav skulle bli til ein faktisk plan, stogga maskineriet. Diagrammet sa kva som høyrde saman. Det sa ikkje korleis det skulle sjå ut. Det siste spranget, frå struktur av krav til verkeleg rom, måtte Alexander ta med same auge og same hand som ein arkitekt utan noko diagram. Den vakre nedbrytinga leverte honom attende til seg sjølv ved terskelen til det einaste som var design. Då han seinare arbeidde med Bay Area Rapid Transit i San Francisco på slutten av sekstitalet, var lærdomen den same, gjenteken prosjekt etter prosjekt: metoden fjerna ikkje dommen, han utsette honom.
\end{kasus}

Her ligg det presise skilet frå ein vanleg falsifisert hypotese. Når ein hypotese i eit modent fag fell, fell han \emph{mot} fundamentet og gjer det sterkare: eksperimentet motseier teorien, men målemetoden, logikken og omgrepa står att og lèt deg formulere ein betre teori. Tapet er produktivt fordi noko under ber tapet. Då metoderørsla fall, stod det ingen ramme att til å bere fallet. Det var ikkje ein bestemt metode som vart motbevist medan ramma heldt; det var sjølve tanken om at design \emph{kunne} ha ein eksplisitt metode som gav etter, fordi det aldri fanst ein teori om form å sikte mot. Difor akkumulerte faget ingenting. Kjemien gjekk frå flogiston til oksygen; design gjekk frå metode til avvising av metode og deretter attende til intuisjonen det starta med, ein heil rundtur som enda der han byrja, no berre meir sjølvmedviten. Eit fag som kan gjere ein slik rundtur fortel deg at det aldri fanst eit fast punkt å bevege seg vekk frå.

Det Alexander gjorde etterpå, er sjølv eit vitnemål. I staden for ein metode skreiv han, saman med andre, \textit{A Pattern Language} (1977): ikkje ein prosedyre for å kome fram til form, men ein katalog over 253 mønster som hadde vist seg å verke, kvart med ei grunngjeving\autocite{alexander1977pattern}. Han gav opp draumen om ein metode som kunne generere det gode, og skreiv i staden ned det gode direkte. Når metoden sviktar fordi fundamentet manglar, kan ein anten lyge om at metoden held, eller gå attende til å katalogisere det ein likar. Faget gjorde for det meste det fyrste. Alexander gjorde det andre, og var iallfall ærleg om kva han då sat med: ein katalog, ikkje ein teori.

\section{Forma som kunne reknast}

Same året som Alexander bad alle gløyme heile greia, 1971, og året etter, la to unge forskarar fram nett det rørsla hadde sukka etter i eit tiår utan å greie å byggje sjølv: ein formell, generativ, utreknbar teori om form. George Stiny og James Gips kalla han \emph{formgrammatikk}, shape grammar\autocite{stiny1972}. Tanken er enkel. Ein formgrammatikk er eit lite sett reglar som opererer direkte på former, ikkje på symbol som står for former: ein startfigur, og reglar av typen «når du ser denne forma, kan du byte henne ut med denne». Køyrer du reglane, genererer dei eit heilt formrom, og rommet er ikkje ein vag metafor, men ei presist definert mengd ein kan rekne i, telje i og prøve påstandar mot.

Dette treffer holet vi nett har grave fram. Metoden fall fordi det ikkje fanst nokon teori om form å sikte mot; stega førte heilt til kanten av valet og slapp taket der. Formgrammatikken seier ikkje kva som er vakkert, men han seier presist kva som kan setjast saman av kva, kva som ligg innanfor og utanfor eit gjeve formspråk, og korleis ein stil kan skrivast som ein endeleg regelmengd snarare enn som ei kjensle i handa. Stiny og William Mitchell skreiv Andrea Palladios villaer som ein grammatikk: eit regelsett som genererer planane til dei verkelege villaene, og uendeleg mange nye av same slag, som Palladio kunne ha teikna men aldri gjorde\autocite{stinymitchell1978}. Terry Knight førte apparatet vidare til det faget hadde gjeve opp å gjere strengt, sjølve stilendringa, skriven som transformasjonar av ein grammatikk\autocite{knight1994}. Frå 1972 til 1978 til 1994 går det ei kumulativ line med eige tidsskriftstoff, eigne konferansar og resultat som byggjer på kvarandre.

Lina vaks ikkje i designskulane. Han voks i randsonene til arkitektur og databehandling, i computational design, i nokre tekniske miljø, og let seg aldri innlemme i den allmenne designutdanninga som det fundamentet han kunne ha vore. Den jamne industri-, interaksjons- eller tenestedesignaren går ut av eit femårig løp utan å ha høyrt namnet. Faget kan ikkje eingong seie at gåva var framand, for ho kom ikkje frå ein økonom utanfrå, men frå folk som arbeidde med form, om form, på formas eigne premiss. Det lét henne liggje fordi ho kravde matematikk, og matematikk var prisen faget heile tida nekta å betale.

\vignett

Metoderørsla var ein av fleire framstøytar i samtida for å gjere det menneskeskapte til gjenstand for ein eigen vitskap. Herbert Simon argumenterte parallelt for at design kunne bli ein streng vitskap om det kunstige; han får eit eige kapittel, så her er det nok å sjå at han delte rørsla si vone om at det rette resultatet ville følgje om ein berre fekk prosessen rasjonell nok. Det er den vona kapitlet har vist er tom: ein rasjonell prosess kan berre rekne seg fram til det rette når det finst eit mål utanfor prosessen som seier kva rett er. Bruce Archer prøvde å berge ambisjonen ved å heve han eit nivå. Han slutta ikkje at design ikkje kunne vere stringent, men at design måtte konstituerast som ein eigen \emph{disiplin} med ein kunnskapsform skild frå både vitskap og humaniora\autocite{archer1979discipline}. I det grepet vert mangelen sjølv omdøypt til ein eigenart: designen har ikkje den vitskaplege metoden sitt grunnlag, \emph{difor} er han ein annan type kunnskap med sin eigen rett. Det er faget sin mest brukte manøver, og her dukkar han opp fyrste gong som svar på metodekrisa. Horst Rittel, sjølv ein gong metodemann frå Ulm, drog den hardaste slutninga: saman med Melvin Webber formulerte han i 1973 at planleggings- og designproblem er \emph{wicked problems}, vrange problem utan endeleg formulering, utan rett løysing, utan ein test som seier om dei er løyste\autocite{rittel1973}. Som realistisk skildring er det sant. Men for faget vart det ei trøyst: ein grunn til å slutte å leite. Om problema er vrange av natur, treng ingen skemmast over at metoden aldri vart funnen. Den vesle vakre tanken kom til å bere ei tung last i tiåra etter, og det er den lasta neste kapittel skal vege.

\laeringsmal{\begin{itemize}
\item Gjere greie for designmetoderørsla (konferansen 1962, Jones, Alexander) som forsøk på å gje design ein eksplisitt, rasjonell metode.
\item Forklare kvifor grunnleggjarane sjølve forkasta rørsla, og kva Alexanders avvising i \textit{DMG Newsletter} frå 1971 viser om det tomme valsteget.
\item Skjøne kvifor formgrammatikken (Stiny, Mitchell, Knight) var den strenge teorien om form faget aldri tok imot.
\item Kjenne att «Archer-manøveren»: å omdøype eit fråvær (manglande metode) til ein eigenart (ein eigen kunnskapsform), og knyte metodekrisa til Rittels vrange problem.
\end{itemize}}

\ovingar{\begin{enumerate}
\item Alexander sa at han ikkje ville snakke om noko kalla «metodologi». Var avvisinga ei innrømming av nederlag, eller ei sunn korreksjon? Argumenter for begge sider før du konkluderer.
\item «Å gjere eit fråvær om til ein eigenart.» Finn eit anna fag eller felt som gjer same manøveren. Held han betre der enn i design?
\item Skisser eit lite eksperiment som kunne avgjere om ein føreslått designmetode faktisk gjev betre utfall enn inga metode. Kvifor er slike eksperiment sjeldne i faget?
\end{enumerate}}

\vidarelesing{\fullcite{jones1970methods}; \fullcite{alexander1964notes}; \fullcite{alexander1971state}; \fullcite{alexander1977pattern}; \fullcite{cross1984developments}; \fullcite{archer1979discipline}; \fullcite{rittel1973}}
